\scanpage{45}

On pourra donc poser

\[
\begin{aligned}
\unsure{y|\eta}_s&=\xi_s T_\nu^{\delta_\nu}T_{\unsure{\nu-1|\nu+1}}^{\delta_{\nu+1}}\cdots T_{r-1}^{\delta_{r-1}}(T_r,\ldots T_{s-1})^{p-1}\\
&\text{avec }\xi_s\in M
\end{aligned}
\]

On aura

\[
y_s T_s^{p-1}=\xi_s T_\nu^{\delta_\nu}\cdots T_{r-1}^{\delta_{r-1}}(T_r\cdots T_s)^{p-1}
\]

La relation $(R_s)$ s'écrit donc, comme

\[
p\unsure{\gamma|\nu}'_s y_{s+1}=p\gamma'_s\xi_{s+1}T_\nu^{\delta_\nu}\cdots T_{r-1}^{\delta_{r-1}}(T_r\cdots T_s)^{p-1}
\]

\[
\boxed{\xi_s=p\gamma'_s\xi_{s+1}}
\]

On aura donc

\[
\xi_s=p^h\gamma'_s\cdots\gamma'_{s+h-1}\,\xi_{s+h}\qquad\forall\,h\geq1
\]

\[
\text{donc}\qquad \xi_s\in\deleted{\unsure{p}}\bigcap_h(p^h M)
\]

i.e. les $\xi_s$ appartiennent à la partie divisible de $\underline{M}$.

\[
\leadsto
\]

\uline{Prop} Soit $k$ un anneau comm, $\deleted{\unsure{B}}S_k=k\{T\}$, $M,N$ deux $k$-modules, \deleted{avec} $d>0$, $N_0^{\unsure{*|\natural}}=\bigcap_{n\in\mathbb{N}^{*}}nN$ le \deleted{partie} \inserted{plus grand} sous-module divisible de $N$, $\alpha\geq d$. Alors si

$u\in\mathrm{Hom}^{-d}(\tau_\alpha M_{S_k},N_{S_k})$ (\deleted{qui} \inserted{lequel module} est indépendant \unsure{mod restriction|mod restrictions} du choix de $\alpha$), $u$ se factorise par $N_{0\,S_k}$, i.e.

\[
\mathrm{Hom}^{-d}(\tau_\alpha M_{S_k},N_{0\,S_k})\overset{\sim}{\unsure{\longleftrightarrow|\longrightarrow}}\mathrm{Hom}^{-d}(\tau_\alpha M_{S_k},N_{0\,S_k})
\]

\uline{Cor} Si $N_0=0$, p. ex. $k$ de torsion, alors $\mathrm{Hom}^{-d}(\tau_\alpha M_{S_k},N_{S_k})=0$

% archive sheet number: 44
